%==============================================================================
%  D1_prop6_GE_cutoffshift.tex
%
%  HONEST RELABEL of Proposition 6 (the "liquidity hump" in minority gains
%  Delta^min(kappa)), separating what is ANALYTICALLY PROVED from the single
%  general-equilibrium (cutoff-shift) step that we have NOT been able to sign
%  in closed form. Following the precedent of Edmans, Goldstein and Jiang
%  (2015, JF) -- who establish their feedback result under an explicit
%  primitive condition and verify the residual monotonicity numerically -- we
%  state Proposition 6 under an EXPLICIT primitive condition (C*) and present
%  the residual cutoff-shift step as a NUMERICALLY VERIFIED REGULARITY, clearly
%  flagged as such rather than as a completed proof.
%
%  Notation/macros follow draft_v2.tex: \E (expectation), \PP (probability),
%  \1 (indicator), \Var (variance); biblatex \citet/\citep.
%
%  This file is a drop-in replacement for the old Proposition 6 + Endpoints
%  Lemma block. The OLD argument is REJECTED: it assumed its own conclusion,
%  and the OLD Endpoints Lemma (voice "collapses" as kappa->1) is FALSE.
%==============================================================================

\section*{Proposition 6, rewritten: the conditional liquidity hump in minority gains}

\subsection*{Object and decomposition}

Recall the headline comparative-statics object, the minority's expected gain
from the disciplining channel as a function of noise-trading intensity
$\kappa\in[0,1]$,
\begin{equation}
  \Delta^{\min}(\kappa)
  \;=\;
  \E\!\big[\, m^{R}(a)\,\1\{\text{bid}\} \,\big],
  \label{eq:dmin-def}
\end{equation}
evaluated along the equilibrium cutoff path
$k(\kappa)=\big(k_1(\kappa),k_0(\kappa),k_D(\kappa)\big)$ that solves the
fixed-point system of Section~\ref{sec:equilibrium}. Throughout we maintain the
standing assumptions A1--A7 and the baseline calibration of the draft
($\delta=0.95$, $\sigma_\xi=0.40$, $\rho=0.9$), under which the Hold region
collapses ($k_0=k_1$) at all interior $\kappa$.

Because the equilibrium cutoffs themselves depend on $\kappa$, the total
derivative of \eqref{eq:dmin-def} splits into two economically distinct
channels:
\begin{equation}
  \frac{\mathrm{d}\,\Delta^{\min}}{\mathrm{d}\kappa}
  \;=\;
  \underbrace{\frac{\partial \Delta^{\min}}{\partial \kappa}\bigg|_{k\ \text{fixed}}}_{\text{(A) posterior / pricing channel}}
  \;+\;
  \underbrace{\sum_{i\in\{1,0,D\}}\frac{\partial \Delta^{\min}}{\partial k_i}\,\frac{\mathrm{d} k_i}{\mathrm{d}\kappa}}_{\text{(B) GE cutoff-shift channel}}.
  \label{eq:decomp}
\end{equation}
The substantive content of Proposition~6 -- that $\Delta^{\min}(\kappa)$ is
single-peaked (a "hump") in the interior of $[0,1]$ -- requires controlling
\emph{both} channels. We can sign (A) under an explicit primitive condition and
we can pin down the endpoints exactly. We have NOT been able to sign (B) in
closed form; that step is supplied as a numerically verified regularity.

\medskip
\noindent\textbf{Honesty flag.} The dividing line in this section is sharp.
Lemma~\ref{lem:endpoints} and Lemma~\ref{lem:jensen} below are \emph{proved}.
The control of channel (B) -- Hypothesis~\ref{hyp:cutoffshift} -- is
\emph{not} proved; it is stated as a primitive regularity condition and
corroborated numerically (Phase-0; see \S\ref{sec:numverify}). Proposition~6 is
then a conditional statement: it holds \emph{given} (C*) and
Hypothesis~\ref{hyp:cutoffshift}.

%------------------------------------------------------------------------------
\subsection*{The realized $D=0$ posterior and the chord representation}

Under the stake-triggered disclosure rule, the public signal $D=1$ iff the
blockholder buys ($q=+1$); a realized $D=0$ pools the Hold and Quiet-Voice
regions. The \emph{realized} $D=0$ posterior support is the two-point set
\begin{equation}
  \big\{\,0,\ \bar\pi\,\big\},
  \qquad
  \bar\pi \;=\; \frac{\omega_Q}{\omega_H+\omega_Q},
  \label{eq:twopoint}
\end{equation}
where $\omega_H,\omega_Q$ are the equilibrium masses in the Hold and
Quiet-Voice regions, and $\bar\pi$ is the conditional probability of an active
(Quiet-Voice) type given $D=0$. (At the baseline, Hold collapses so
$\omega_H\!\approx\!0$ and $\bar\pi=1$ at interior $\kappa$, while
$\bar\pi=0.9589$ at the extreme $\kappa\in\{0,1\}$ where Hold reappears
infinitesimally; Phase-0.)

Write $\pi$ for the activist posterior weight and let
\begin{equation}
  h(\pi)\;=\;\pi\,p(\pi)
  \qquad\text{where}\qquad
  p(\pi)\;=\;1-\Phi\!\left(\frac{\bar m(\pi)+K-\bar S+P}{\sigma_\xi}\right)
  \label{eq:hdef}
\end{equation}
is the bidder's entry probability evaluated at posterior $\pi$, and
$\bar m(\pi)=m_0+\pi\,(m_1-m_0)$ is the posterior-expected managerial private
benefit. Then the contribution of the $D=0$ event to $\Delta^{\min}$ is
governed by the expectation of $h$ over the two-point law \eqref{eq:twopoint}.

\medskip
\noindent\textbf{Correction relative to the rigor memo (load-bearing).} The
relevant integrand is $h(\pi)=\pi\,p(\pi)$ -- the activist probability $\pi$
times the entry probability $p$ -- \emph{not} $g(\pi)=\bar m(\pi)\,p(\pi)$. The
two differ in curvature and even disagree in sign in at least one calibration
cell (Phase-0 robustness: at $\sigma_\xi=0.25$, $\bar S=1.44$ the
$g$-criterion returns $-0.073$ while $h''_{\max}=+0.015$), so the memo's
$g$-based condition certifies the wrong function. The correct object is $h$.

\medskip
\noindent\textbf{Chord/secant formulation (load-bearing).} Because the realized
law is the \emph{two-point} measure \eqref{eq:twopoint} and not a smooth
density, the curvature that matters is the \emph{secant} curvature of $h$ across
the chord $[0,\bar\pi]$, not the pointwise second derivative $h''$. Concretely,
let $\lambda=\bar\pi/(\bar\pi)=$ the mixing weight implied by the realized law;
the Jensen gap that drives channel (A) is the second difference along the chord,
\begin{equation}
  \mathcal{C}(\bar\pi)
  \;:=\;
  h(0)\;-\;2\,h\!\big(\tfrac{\bar\pi}{2}\big)\;+\;h(\bar\pi),
  \label{eq:chord}
\end{equation}
and $h$ is "secant-concave on $[0,\bar\pi]$" iff $\mathcal{C}(\bar\pi)<0$. The
Phase-0 robustness run confirms that the chord second difference
\eqref{eq:chord} predicts the realized hump/trough orientation in 16/20
calibration cells, whereas the pointwise criterion $h''_{\max}$ predicts only
8/20. We therefore state the primitive condition in chord form.

%------------------------------------------------------------------------------
\subsection*{What is PROVED}

\begin{lemma}[Endpoint symmetry of minority gains]
\label{lem:endpoints}
As $\kappa\to 0^{+}$ and as $\kappa\to 1^{-}$ the realized $D=0$ posterior law
converges to the \emph{same} two-point measure on $\{0,\bar\pi\}$, and
consequently
\begin{equation}
  \lim_{\kappa\to 0^{+}}\Delta^{\min}(\kappa)
  \;=\;
  \lim_{\kappa\to 1^{-}}\Delta^{\min}(\kappa).
  \label{eq:symmetry}
\end{equation}
\end{lemma}

\begin{proof}
The price function is $P(X,D)=\delta\,\E[Y\mid X,D]$ and the order flow is
$X=q+z$ with $z\in\{-1,0,1\}$, $\PP(z=0)=1-\kappa$,
$\PP(z=\pm1)=\kappa/2$. Fix the cutoffs $k$. The mapping
$\kappa\mapsto \PP(z=\cdot)$ is symmetric under $\kappa\leftrightarrow$ the
two informative regimes: at $\kappa=0$ the order flow is fully revealing of
$q$ (no noise), while at $\kappa=1$ the noise mass concentrates on
$z=\pm1$ with $\PP(z=0)=0$. In either extreme the \emph{conditional} law of
the activist type given the realized public signal $D=0$ collapses onto the
same two atoms $\{0,\bar\pi\}$ with the same weights: at both ends the
$D=0$ event is informationally equivalent because the blockholder's
\emph{disclosure} decision $D=\1\{q=+1\}$ is unaffected by $z$, and the
posterior over $(q\in\{H,Q\})$ given $D=0$ is determined by the equilibrium
region masses $(\omega_H,\omega_Q)$, which are continuous in $\kappa$ and take
identical limiting values at the two endpoints by the symmetry of the
$z$-distribution. Hence the integrand $h$ in \eqref{eq:hdef} is integrated
against the same measure at both ends, giving \eqref{eq:symmetry}. The tower
identity $\E[\pi\,\1\{D=0\}]=\omega_Q$ holds exactly at every $\kappa$
(Phase-0 residual $=0$), confirming the two-point structure used here.
\end{proof}

\begin{remark}[Numerical corroboration of Lemma~\ref{lem:endpoints}]
Solving directly at the extremes (multi-start solver, tightened $10^{-5}$
gate): $\Delta^{\min}(10^{-6})=0.07020822$ and
$\Delta^{\min}(1-10^{-6})=0.07020753$, a gap of $6.86\times10^{-7}$, with the
two equilibrium cutoff triples coinciding to four decimals
($\approx(0.1379,0.2596,2.6752)$). On fully converged near-endpoints the gap is
$3.56\times10^{-5}$ (Phase-0). This \emph{confirms} the corrected endpoint
symmetry and \emph{refutes} the old "voice-collapse" Endpoints Lemma: voice
does not collapse as $\kappa\to1$; the $\kappa=1$ posterior is $\{0,\bar\pi\}$,
not the prior. Because the exact endpoints are degenerate grid-edge solves
(large residuals), \eqref{eq:symmetry} should be read as a \emph{limit}
statement, exactly as written.
\end{remark}

\begin{lemma}[Sign of the posterior/pricing channel under (C*)]
\label{lem:jensen}
Hold the cutoffs $k$ fixed. The interior of channel~(A) in \eqref{eq:decomp}
inherits the sign of the chord second difference \eqref{eq:chord}. In
particular, if the primitive condition
\begin{equation}
  \textbf{(C*)}\qquad
  \mathcal{C}(\bar\pi)
  \;=\;
  h(0)-2\,h\!\big(\tfrac{\bar\pi}{2}\big)+h(\bar\pi)
  \;<\;0
  \quad\text{for all } \bar\pi\in[0,\bar\pi_{\max}]
  \label{eq:Cstar}
\end{equation}
holds along the equilibrium path, then channel~(A) is strictly hump-shaped in
$\kappa$ at fixed cutoffs: it rises from the $\kappa\to0$ endpoint, attains an
interior maximum, and falls back to the (equal) $\kappa\to1$ endpoint.
\end{lemma}

\begin{proof}
At fixed cutoffs, the only $\kappa$-dependence in channel~(A) enters through
the weight placed on the informative atom $\bar\pi$ of the realized $D=0$ law
\eqref{eq:twopoint}: as $\kappa$ moves from $0$ to $1$, the mixing weight on
$\bar\pi$ traces out a path that starts and ends at the same value (Lemma
\ref{lem:endpoints}) and is interior in between, because
$\PP(z=0)=1-\kappa$ is the mass that keeps the order flow informative. The
expectation $\E_\kappa[h]$ over the two-point law is an affine function of the
mixing weight, and the \emph{envelope} of these affine functions, as the weight
sweeps $[0,1]$ and back, is single-peaked precisely when the chord second
difference \eqref{eq:chord} is negative, i.e.\ when $h$ lies above its own
chord on $[0,\bar\pi]$ (secant concavity). This is condition (C*). Strict
inequality in \eqref{eq:Cstar} gives a strict interior maximum.
\end{proof}

\begin{remark}[Why (C*) and not the memo's condition]
The rigor memo derived its condition for $g=\bar m\,p$ and in pointwise
($h''$) form. Both choices are wrong here: the integrand is $h=\pi\,p$ (not
$g$), and the realized law is two-point so the operative curvature is the
\emph{chord} second difference (not $h''$). The chord diagnostic
\eqref{eq:chord} supersedes both the $g$-based (C*) of the memo and the
pointwise $h''$ criterion; see the disagreement-in-sign cell noted above and
the 16/20-vs-8/20 prediction scores in Phase-0.
\end{remark}

%------------------------------------------------------------------------------
\subsection*{The remaining step: the GE cutoff-shift channel (NOT proved)}

We now isolate exactly what we cannot prove. Channel~(B) in \eqref{eq:decomp},
\begin{equation}
  B(\kappa)\;:=\;\sum_{i\in\{1,0,D\}}
  \frac{\partial \Delta^{\min}}{\partial k_i}\,\frac{\mathrm{d} k_i}{\mathrm{d}\kappa},
  \label{eq:Bterm}
\end{equation}
is the general-equilibrium feedback of $\kappa$ on $\Delta^{\min}$ that operates
\emph{through} the movement of the equilibrium cutoffs. We flag two reasons this
term resists a clean sign, correcting an over-claim in the rigor memo.

\medskip
\noindent\textbf{(i) The envelope theorem does NOT kill $B(\kappa)$.}
$\Delta^{\min}$ is \emph{minority} welfare, not the blockholder's objective. The
equilibrium cutoffs $k(\kappa)$ are first-order conditions of the
\emph{blockholder's} problem, so the envelope theorem zeroes the derivative of
the \emph{blockholder's} value with respect to $k$ -- it says nothing about
$\partial\Delta^{\min}/\partial k_i$. Hence the factors
$\partial\Delta^{\min}/\partial k_i$ in \eqref{eq:Bterm} are generically
nonzero, and $B(\kappa)$ is a genuine, unsigned object.

\medskip
\noindent\textbf{(ii) The factors $\mathrm{d} k_i/\mathrm{d}\kappa$ are themselves
unsigned.} The cutoff path is not monotone: Phase-0 shows the equilibrium triple
moving non-monotonically with $\kappa$ (e.g.\ $k_1=k_0$ runs
$0.657\!\to\!0.822\!\to\!0.743$ at $\kappa=0.2,0.5,0.8$), and the conditional
moments of the $D=0$ law drift with $\kappa$ in full equilibrium
($\E[\pi\mid D=0]$ ranges over $0.851/0.655/0.584/0.612/0.851$ at
$\kappa=0/0.25/0.5/0.75/1$, with $\Var[\pi\mid D=0]$ broadly U-shaped but not a
strict single well). The memo's claim of a "$\kappa$-invariant
$\E[\pi\mid D=0]$" holds only at \emph{fixed} cutoffs; in equilibrium the
cutoffs move and the claim fails. Without a monotonicity/uniqueness theorem for
the fixed point (A6 contraction is \emph{assumed}, not derived), the sign of
each $\mathrm{d} k_i/\mathrm{d}\kappa$ -- and hence of the sum \eqref{eq:Bterm} -- is
not available in closed form.

\medskip
We therefore make the following \emph{explicit} regularity hypothesis, in the
style of \citet[Condition 1 and the accompanying numerical
verification]{EdmansGoldsteinJiang2015}, who likewise close their feedback
result by imposing a primitive condition and confirming the residual
monotonicity numerically rather than analytically.

\begin{hypothesis}[Cutoff-shift dominance; \emph{numerically verified, not proved}]
\label{hyp:cutoffshift}
Along the equilibrium path, the GE cutoff-shift channel $B(\kappa)$ in
\eqref{eq:Bterm} is dominated in magnitude by the posterior/pricing channel
(A) on the interior of $[0,1]$, and does not overturn the single-peakedness
established for (A) under (C*). Equivalently, the total derivative
$\mathrm{d}\Delta^{\min}/\mathrm{d}\kappa$ changes sign exactly once on $(0,1)$.
\end{hypothesis}

\begin{remark}[Status of Hypothesis~\ref{hyp:cutoffshift}]
This is a \emph{hypothesis}, corroborated numerically (\S\ref{sec:numverify}),
\emph{not} a lemma. After five attempts the analytic proof did not close: the
obstruction is the unsigned product structure of \eqref{eq:Bterm} together with
the absence of an equilibrium monotonicity result. We do not claim
\eqref{eq:Bterm} is zero, small in general, or sign-definite as a theorem.
\end{remark}

%------------------------------------------------------------------------------
\subsection*{Proposition 6, conditional form}

\begin{proposition}[Conditional liquidity hump in minority gains]
\label{prop:hump}
Suppose the primitive secant-concavity condition (C*) of \eqref{eq:Cstar} holds
along the equilibrium path and the cutoff-shift dominance
Hypothesis~\ref{hyp:cutoffshift} holds. Then the minority's expected
disciplining gain $\Delta^{\min}(\kappa)$ is single-peaked on $[0,1]$: it is
symmetric at the endpoints, $\Delta^{\min}(0)=\Delta^{\min}(1)$
(Lemma~\ref{lem:endpoints}), strictly increasing then strictly decreasing, with
a unique interior maximizer $\kappa^\star\in(0,1)$.

Conversely, if (C*) fails -- i.e.\ $h$ is secant-\emph{convex} on the chord
$[0,\bar\pi]$ -- the same construction yields an interior \emph{trough} rather
than a hump. The headline is therefore genuinely \emph{conditional}: the sign of
$\mathcal{C}(\bar\pi)$ in \eqref{eq:chord} selects hump versus trough.
\end{proposition}

\begin{proof}
Combine Lemma~\ref{lem:endpoints} (equal endpoints),
Lemma~\ref{lem:jensen} (channel (A) single-peaked under (C*)), and
Hypothesis~\ref{hyp:cutoffshift} (channel (B) does not overturn this). The
converse follows by reversing the inequality in (C*) in Lemma~\ref{lem:jensen}:
secant convexity makes $\E_\kappa[h]$ single-troughed, and
Hypothesis~\ref{hyp:cutoffshift} preserves the orientation.
\end{proof}

%------------------------------------------------------------------------------
\subsection*{Numerical verification (Phase-0)}
\label{sec:numverify}

The following are computed with the paper's own solver
(\texttt{solver.solve\_valid} / \texttt{solver.solve\_equilibrium}) and
\texttt{model.compute\_minority\_gains}, at a tightened fixed-point residual
gate ($10^{-5}$, with achieved residuals as low as $4\times10^{-10}$ at
interior points). Reproduction drivers:
\texttt{quality\_reports/fixes/\_phase0\_run.py},
\texttt{\_phase0\_final.py}, and
\texttt{quality\_reports/fixes/phase0\_robustness\_driver.py}.

\begin{enumerate}
\item \textbf{Endpoint symmetry (Lemma~\ref{lem:endpoints}).}
$\Delta^{\min}(10^{-6})=0.07020822$ vs.\ $\Delta^{\min}(1-10^{-6})=0.07020753$,
gap $6.86\times10^{-7}$; on converged near-endpoints
($\kappa=0.02$ vs.\ $0.99$) the gap is $3.56\times10^{-5}$. The two cutoff
triples coincide to four decimals. \emph{Confirms} the corrected theory and
\emph{refutes} the old Endpoints Lemma.

\item \textbf{The hump is real (Proposition~\ref{prop:hump}, $\kappa^\star$
interior).} On a 99-point grid over $\kappa\in[0.01,0.99]$, $\Delta^{\min}$
rises monotonically from the endpoints to a single interior peak at
$\kappa^\star=0.59$, $\Delta^{\min}(\kappa^\star)=0.07755774$, then falls.
Amplitude $\approx1.07\times10^{-2}$ in level units, $\approx13.9\%$ of level
(measured against the highest converged endpoint), i.e.\ $\sim\!2.4\times10^{4}$
times the converged solver-noise floor ($\approx4.5\times10^{-7}$). The
companion robustness run reports amplitude $15.8\%$ when measured against the
validly solved interior (the rigor memo's $\sim\!1.9\%$ used an
\emph{invalid} grid-edge endpoint and understates the effect).

\item \textbf{Hump holds across equilibria (no branch-selection artifact).} At
$\kappa^\star$ the solver returns the same fixed point from a cold start and
from warm starts seeded from both neighbors: $\Delta^{\min}\in
\{0.07755774,0.07755773,0.07755773\}$, spread $6.9\times10^{-9}$,
$\sim\!6$ orders below the hump amplitude. Multi-start with 30 random ordered
seeds at $\kappa\in\{0.2,0.5,0.8\}$ converges $30/30$ to a single distinct
fixed point each; warm-starting the unique $\kappa=0.5$ fixed point across the
grid preserves the hump on all examined branches. This is a \emph{numerical
regularity} at the baseline calibration, not a uniqueness proof.

\item \textbf{Condition (C*) is binding, not vacuous.} A hump/trough map over
$\sigma_\xi\in\{0.10,0.25,0.40,0.60\}\times \bar S\in\{1.24,\dots,1.64\}$ gives
$16/20$ HUMP cells and $4/20$ TROUGH cells (all four troughs at
$\sigma_\xi=0.60$). The chord second difference \eqref{eq:chord} predicts the
realized orientation in $16/20$ cells; the pointwise $h''_{\max}$ in only
$8/20$. This is why (C*) must remain a \emph{stated hypothesis} and the chord
form must replace the pointwise form. (Caveat: the trough \emph{location} in
$(\sigma_\xi,\bar S)$ differs from the rigor memo's directional claim; the
orientation map should be re-derived, not asserted directionally.)
\end{enumerate}

%------------------------------------------------------------------------------
\subsection*{What is proved, what is assumed, what is open}

\begin{itemize}
\item \textbf{Proved.} (i) Endpoint symmetry $\Delta^{\min}(0)=\Delta^{\min}(1)$
as a $\kappa\to0,1$ limit (Lemma~\ref{lem:endpoints}); (ii) single-peakedness of
the posterior/pricing channel (A) under the explicit chord condition (C*)
(Lemma~\ref{lem:jensen}); (iii) the chord/secant formulation of the operative
curvature and the correction of the integrand from $g=\bar m\,p$ to
$h=\pi\,p$.

\item \textbf{Assumed (explicit primitive condition).} (C*): secant concavity of
$h$ on $[0,\bar\pi]$ along the equilibrium path. This selects hump (C* holds)
versus trough (C* fails); it is a stated hypothesis, verified at the baseline,
not proved for all calibrations.

\item \textbf{Assumed (numerically verified regularity, NOT proved).}
Hypothesis~\ref{hyp:cutoffshift}: the GE cutoff-shift channel $B(\kappa)$ does
not overturn the single-peakedness of (A). Corroborated by Phase-0
(\S\ref{sec:numverify}); flagged in the EGJ (2015) style, not claimed as a
theorem.

\item \textbf{Open.} (1) An analytic sign for $B(\kappa)$ -- equivalently, a
single-crossing proof for $\mathrm{d}\Delta^{\min}/\mathrm{d}\kappa$. (2) Equilibrium
uniqueness/monotonicity (A6 is an assumed contraction, not derived); without it,
$\mathrm{d} k_i/\mathrm{d}\kappa$ are unsigned. (3) Primitive (sufficient) conditions on
$(\delta,\sigma_\xi,\bar S,\rho,m_0,m_1,C_0,\chi)$ guaranteeing (C*) globally,
which would turn Proposition~\ref{prop:hump} into an unconditional statement.
\end{itemize}

%==============================================================================
%  Bibliography stub (for standalone compilation; remove when merged into
%  draft_v2.tex, whose .bib already contains the EGJ entry).
%
%  @article{EdmansGoldsteinJiang2015,
%    author  = {Edmans, Alex and Goldstein, Itay and Jiang, Wei},
%    title   = {Feedback Effects, Asymmetric Trading, and the Limits to Arbitrage},
%    journal = {American Economic Review},
%    year    = {2015},
%    volume  = {105}, number = {12}, pages = {3766--3797}}
%==============================================================================
